% SOURCE_AUTHORITY: sga5_fr_workpass.tex, lines 12792--12834 (LF-normalized unit).
% SOURCE_SHA256: 791F4EFFC5E02832D5D77ED03518C8156D6F07E4C8238B03545DB93D883FBB28
Denotemos por $\tau$ o primeiro membro desta equação, e seja $g\in G$. O elemento $g$ induz evidentemente um endomorfismo do triângulo distinguido
\[
\begin{tikzcd}[column sep=huge,row sep=large]
& \RG_{Y'}(\Lambda_{n,Y'}) \arrow[dl] \arrow[dr] & \\
\RG_{C'}(\Lambda_{n,U',C'}) \arrow[rr] && \RG_{C'}(\Lambda_{n,C'})
\end{tikzcd}
\]
Tem-se, portanto,
\[
\tau(g)=\Tr^*g_{\RG_{C'}(\Lambda_{n,U',C'})}
=\Tr^*g_{\RG_{C'}(\Lambda_{n,C'})}
-\Tr^*g_{\RG_{Y'}(\Lambda_{n,Y'})}.
\]
Para abreviar, escreveu-se $\Tr^*$ em vez de $\Tr^*_{\Lambda_n}$.

Ora, aplicando uma fórmula de Lefschetz trivial ao espaço discreto $Y'$ sobre o qual opera $g$, tem-se:
\[
\begin{aligned}
\Tr^*g_{\RG_{Y'}(\Lambda_{n,Y'})}
&=\text{número de pontos fixos da restrição } g_{Y'} \text{ de } g \text{ a } Y'\\
&=\sum_{\substack{y'\in Y'\\gy'=y'}}1.
\end{aligned}
\]
Por outro lado, a fórmula de Lefschetz (III 4.7) para a ação de $g$ sobre $C'$ dá, se $g\ne e$,
\[
\Tr^*g_{\RG_{C'}(\Lambda_{n,C'})}
=
\sum_{y'\in Y'} v_{y'}(g),
\]
a soma das multiplicidades dos pontos fixos da aplicação $g_{C'}:C'\to C'$, e tem-se
\[
\Tr^*g_{\RG_{C'}(\Lambda_{n,C'})}=\EP(C')\cdot1_\Lambda
\quad\text{para }g=e,
\]
graças ao cálculo explícito dos $H^1(C',\Lambda_{n,C'})$, mediante (SGA 4 IX 47) e Künneth (SGA 4 XVII). Obtém-se, pois, a fórmula seguinte para $\tau$:
\[
\tau(g)=
\begin{cases}
\displaystyle\sum_{\substack{y'\in Y'\\g(y')=y'}}\bigl(v_{y'}(g)-1\bigr),& g\ne e,\\[0.8em]
\EP(C')-\operatorname{card}(Y'),& g=e.
\end{cases}
\tag{5.2.2}
\]
